% ======================================================================
%  环论研究论文骨架  (Ring Theory)
%  编译方式: xelatex -> bibtex -> xelatex -> xelatex
%    (xelatex ring-theory-paper && bibtex ring-theory-paper &&
%     xelatex ring-theory-paper && xelatex ring-theory-paper)
%  说明: 所有【TODO】均为待填充占位, 填充后删除注释即可。
% ======================================================================
\documentclass[11pt]{ctexart}

% ---------- 页面与编码 ----------
\usepackage[a4paper,margin=2.5cm]{geometry}
\usepackage[T1]{fontenc}

% ---------- 数学 ----------
\usepackage{amsmath,amssymb,amsthm}
\usepackage{mathtools}
\usepackage{bm}

% ---------- 表格 / 图 ----------
\usepackage{booktabs}
\usepackage{graphicx}
\usepackage{caption}
\usepackage{subcaption}

% ---------- 排版 ----------
\usepackage{enumitem}
\usepackage{multirow}
\usepackage{array}

% ---------- 参考文献 ----------
\usepackage[numbers,sort&compress]{natbib}

% ---------- 定理环境 (环论常用) ----------
\theoremstyle{plain}
\newtheorem{theorem}{定理}[section]
\newtheorem{lemma}[theorem]{引理}
\newtheorem{proposition}[theorem]{命题}
\newtheorem{corollary}[theorem]{推论}
\theoremstyle{definition}
\newtheorem{definition}[theorem]{定义}
\newtheorem{example}[theorem]{例}
\theoremstyle{remark}
\newtheorem{remark}[theorem]{注}

% ---------- 常用数学记号 (环论) ----------
\newcommand{\Z}{\mathbb{Z}}
\newcommand{\Q}{\mathbb{Q}}
\newcommand{\R}{\mathbb{R}}
\newcommand{\C}{\mathbb{C}}
\newcommand{\F}{\mathbb{F}}
\newcommand{\Ring}{\mathcal{R}}
\newcommand{\id}{\mathrm{id}}
\newcommand{\spec}{\mathrm{Spec}}
\newcommand{\Hom}{\mathrm{Hom}}
\newcommand{\End}{\mathrm{End}}

% ---------- 颜色 (hyperref 需要) ----------
\usepackage{xcolor}

% ---------- hyperref (必须最后加载) ----------
\usepackage[colorlinks=true,
            linkcolor=blue,
            citecolor=green!50!black,
            urlcolor=cyan,
            bookmarksnumbered=true]{hyperref}

% ======================================================================
\begin{document}

% ---------- 标题与作者 ----------
\title{环论中\ 【TODO: 补全论文主题，如“××环的结构刻画”】}
\author{【TODO: 作者姓名】\thanks{【TODO: 单位、邮箱、资助信息】}}
\date{\today}
\maketitle

% ---------- 摘要 ----------
\begin{abstract}
【TODO: 按“五句法则”写中文摘要：(1) 背景一句；(2) 问题一句；(3)
方法一至两句；(4) 主要结果/实验一至两句；(5) 意义一句。】

\noindent\textbf{关键词：}【TODO: 3--6 个，如】环论；交换环；诺特环；同调代数；结构定理
\end{abstract}

% ======================================================================
\section{引言}
\label{sec:intro}

\subsection{研究背景}
\label{sec:background}
【TODO: 交代研究领域。示例句式：】
环作为带乘法与加法的代数结构，是抽象代数的核心对象之一，在代数几何、
数论与表示论中均扮演基础性角色。交换环理论起源于数论中理想与整除性
的抽象化，\cite{placeholder-atiyah} 与 \cite{placeholder-eisenbud}
建立了现代交换代数的基本框架。

近年来，\textbf{【TODO: 具体子方向，如“非交换环的××性质”】} 引起了广泛关注，
相关研究在 \textbf{【TODO: 应用领域，如“编码理论 / 密码学 / 量子信息”】}
中具有潜在应用。然而，关于 \textbf{【TODO: 具体对象】} 的 \textbf{【TODO: 具体性质】}，
目前的研究仍然 \textbf{【TODO: 不充分 / 缺少统一框架 / 计算复杂度过高】}。

\subsection{问题陈述}
\label{sec:problem}
本文研究以下问题：

\begin{enumerate}[label=(\arabic*)]
  \item 【TODO: 问题 1，如“刻画满足××条件的环的结构”】
  \item 【TODO: 问题 2，如“给出××性质的充要判据”】
  \item 【TODO: 问题 3，如“设计判定××的算法并分析复杂度”】
\end{enumerate}

\subsection{本文贡献}
\label{sec:contrib}
本文的主要贡献概括如下：

\begin{itemize}
  \item \textbf{【贡献 1】}【TODO: 如：证明了××环同构于××，填补了……的空白】；
  \item \textbf{【贡献 2】}【TODO: 如：给出了××的显式刻画，改进了已有结果
        \cite{placeholder-ref} 中的界】；
  \item \textbf{【贡献 3】}【TODO: 如：提出了××算法，复杂度为 $O(n^2)$，
        并通过实验验证了其效率】。
\end{itemize}

\subsection{组织与记号}
本文余下部分安排如下：第~\ref{sec:related}~节回顾相关工作；
第~\ref{sec:method}~节给出本文方法；第~\ref{sec:experiment}~节描述实验设置与结果；
第~\ref{sec:conclusion}~节总结全文并讨论未来方向。
未加说明的记号沿用 \cite{placeholder-atiyah}：环均指含幺结合环，
理想均指双边理想；$\spec(R)$ 表示环 $R$ 的素谱。

% ======================================================================
\section{相关工作}
\label{sec:related}
本节从以下四个主题梳理与本文相关的已有工作。

\subsection{主题一：\ 【TODO: 如“交换环的结构理论”】}
\label{sec:rel1}
【TODO: 概述该主题的代表性工作与主要结论，并说明与本文的关联。
示例句式：】Atiyah--MacDonald \cite{placeholder-atiyah} 系统论述了交换环的
诺特性、整性与维数理论；【TODO: 补充 2--3 篇关键文献】……
这些工作为本研究提供了理论基础，但它们 \textbf{【TODO: 未覆盖……】}。

\subsection{主题二：\ 【TODO: 如“非交换环与环的推广”】}
\label{sec:rel2}
【TODO: 概述第二主题，重点说明与本文方法的差异。示例句式：】
在非交换情形，\cite{placeholder-ref} 将经典结论推广到 \textbf{【TODO: 具体类别】}……
本文与之不同之处在于 \textbf{【TODO: 差异化定位】}。

\subsection{主题三：\ 【TODO: 如“同调维数与导出范畴方法”】}
\label{sec:rel3}
【TODO: 概述第三主题。示例句式：】同调方法（投射维数、整体维数等）为环的
结构刻画提供了强大工具，见 \cite{placeholder-weibel}……本文借用了其中的
\textbf{【TODO: 具体工具】}，但处理的对象 \textbf{【TODO: 差异】}。

\subsection{主题四：\ 【TODO: 如“计算与算法方面的工作”】}
\label{sec:rel4}
【TODO: 概述第四主题。示例句式：】在计算方面，【TODO: 算法或软件，如
Gr\"obner 基、GAP、Macaulay2】实现了对环论对象的符号计算……
本文算法与上述工作在 \textbf{【TODO: 输入模型 / 复杂度 / 适用范围】} 上存在差异。

% ======================================================================
\section{方法}
\label{sec:method}
\begin{quote}
\textbf{[占位说明]} 本节为核心部分，当前仅给出章节骨架，内容待补充。
填充时应包含：记号与预备、核心构造/算法、理论分析（定理及其证明）。
\end{quote}

\subsection{记号与预备}
\label{sec:method-prelim}
【TODO: 定义本文所需概念，如】设 $R$ 为含幺结合环。回忆 $R$ 的一个
\textbf{左理想} 是……(见定义~\ref{def:ideal})。

\begin{definition}[理想]
\label{def:ideal}
【TODO: 给出理想的精确定义。】
\end{definition}

\subsection{核心构造}
\label{sec:method-core}
【TODO: 描述本文的核心构造/算法，可分小节。示例结构：】
\begin{enumerate}
  \item \textbf{构造一：}【TODO: 对象 / 映射 / 范畴的构造】；
  \item \textbf{构造二：}【TODO: 算法伪码，用 algorithm 环境补充】。
\end{enumerate}

\subsection{理论分析}
\label{sec:method-analysis}
\begin{theorem}[主定理]
\label{thm:main}
【TODO: 陈述主定理。示例句式：】设 $R$ 满足条件 (C)。则 $R$ 同构于
【TODO: 具体结构】，当且仅当 \dots
\end{theorem}

\begin{proof}
【TODO: 给出证明要点。】
\end{proof}

\begin{corollary}
\label{cor:main}
【TODO: 由主定理推出的直接推论。】
\end{corollary}

\subsection{复杂度分析}
\label{sec:method-complexity}
【TODO: 若含算法，给出时间复杂度与空间复杂度结论，如】
上述算法对 $n$ 阶输入的时间复杂度为 $O(n^2)$，空间复杂度为 $O(n)$，
见引理~\ref{lem:complexity}。

\begin{lemma}
\label{lem:complexity}
【TODO: 复杂度引理。】
\end{lemma}

% ======================================================================
\section{实验设置}
\label{sec:experiment}
\begin{quote}
\textbf{[占位说明]} 本节给出实验部分骨架，含实验环境、数据集、评价指标、
基线与结果展示，内容待填充。
\end{quote}

\subsection{实验环境与实现}
\label{sec:exp-env}
【TODO: 描述软硬件环境，如】实验运行于 Intel Xeon 服务器（2.4~GHz，
128~GB 内存），算法以 Python 3.11 实现，符号计算使用 SageMath 10.x。

\subsection{数据集}
\label{sec:exp-data}
【TODO: 描述测试对象，如】测试环取自 \textbf{【TODO: 来源，如 GAP 库 /
随机生成 / 公开基准】}，共 $N$ 个实例，按 \textbf{【TODO: 规模 / 参数】}
分为三组，详见表~\ref{tab:dataset}。

\begin{table}[htbp]
  \centering
  \caption{数据集统计}
  \label{tab:dataset}
  \begin{tabular}{cccc}
    \toprule
    分组 & 实例数 & 规模范围 & 说明 \\
    \midrule
    A & 【TODO】 & 【TODO】 & 【TODO】 \\
    B & 【TODO】 & 【TODO】 & 【TODO】 \\
    C & 【TODO】 & 【TODO】 & 【TODO】 \\
    \bottomrule
  \end{tabular}
\end{table}

\subsection{评价指标}
\label{sec:exp-metric}
【TODO: 如】采用以下指标：(1) \textbf{正确率}：与人工/权威结果一致的比率；
(2) \textbf{运行时间}（秒）；(3) \textbf{内存峰值}（MB）。

\subsection{基线与对比方法}
\label{sec:exp-baseline}
【TODO: 列出 2--3 个基线，如】本文与以下方法对比：
\begin{itemize}
  \item \textbf{基线 1：}【TODO: 名称与引用 \cite{placeholder-ref}】；
  \item \textbf{基线 2：}【TODO: 名称与引用】；
  \item \textbf{本文方法：}第~\ref{sec:method}~节所提算法。
\end{itemize}

\subsection{结果与分析}
\label{sec:exp-results}
\begin{table}[htbp]
  \centering
  \caption{主要结果}
  \label{tab:results}
  \begin{tabular}{lccc}
    \toprule
    方法 & 正确率(\%) & 时间(s) & 内存(MB) \\
    \midrule
    基线 1 & 【TODO】 & 【TODO】 & 【TODO】 \\
    基线 2 & 【TODO】 & 【TODO】 & 【TODO】 \\
    本文    & 【TODO】 & 【TODO】 & 【TODO】 \\
    \bottomrule
  \end{tabular}
\end{table}

【TODO: 对表~\ref{tab:results} 逐行解读，说明本文方法优于基线的原因。】

% ======================================================================
\section{结论}
\label{sec:conclusion}
本文研究了环论中 \textbf{【TODO: 主题】} 的 \textbf{【TODO: 问题】}，主要成果包括：
(1) 【TODO: 成果一，对应贡献 1】；
(2) 【TODO: 成果二】；
(3) 【TODO: 成果三】。
实验结果表明，本文方法在 \textbf{【TODO: 指标】} 上优于已有基线。

\subsection{局限}
\label{sec:limitation}
本文工作仍存在以下局限：其一，【TODO: 如“仅适用于××条件”】；
其二，【TODO: 如“复杂度随维度指数增长”】；
其三，【TODO: 如“实验规模有限”】。

\subsection{未来工作}
\label{sec:future}
未来工作拟从以下方向展开：(1) 将本文结论推广到 \textbf{【TODO: 更一般情形】}；
(2) 探索 \textbf{【TODO: 应用场景】}；(3) 改进算法的 \textbf{【TODO: 复杂度】}。

% ======================================================================
% 参考文献
% ======================================================================
\bibliographystyle{unsrtnat}
\bibliography{refs}

% ---------- 若暂时不想引用，可注释掉上面两行，用下面内联条目 ----------
% \begin{thebibliography}{9}
%   \bibitem{placeholder-atiyah} M.~F.~Atiyah, I.~G.~MacDonald.
%     \emph{Introduction to Commutative Algebra}. Addison-Wesley, 1969.
%   \bibitem{placeholder-eisenbud} D.~Eisenbud.
%     \emph{Commutative Algebra with a View Toward Algebraic Geometry}.
%     Springer, 1995.
%   \bibitem{placeholder-weibel} C.~A.~Weibel.
%     \emph{An Introduction to Homological Algebra}. Cambridge Univ.\ Press, 1994.
% \end{thebibliography}

\end{document}
